% ----------------------------------------------------------
% Introdução 
% Capítulo sem numeração, mas presente no Sumário
% ----------------------------------------------------------

\chapter*[Introdução]{Introdução}
\addcontentsline{toc}{chapter}{Introdução}

A evolução das telecomunicações espaciais, iniciada na segunda metade do século XX com o lançamento do satélite Sputnik 1 em 1957, transformou a infraestrutura global de dados e monitoramento ambiental \cite{Ford2008}.
Atualmente, o cenário aeroespacial é marcado pela expansão massiva de constelações em órbita terrestre baixa (\textit{Low-Earth Orbit} -- LEO), que operam em altitudes reduzidas e exigem rastreio constante para a recepção de telemetria, imagens de sensoriamento remoto e sinais meteorológicos de alta resolução \cite{app12199505}.

Em aplicações de radiofrequência (RF) envolvendo satélites LEO, a maximização da razão sinal-ruído ($SNR$) e a mitigação de atenuações atmosféricas dependem diretamente do ganho e da diretividade da antena utilizada na estação de solo \cite{hubregt2014}.
Entretanto, antenas de alta diretividade (como arranjos helicoidais ou refletores parabólicos) concentram a energia irradiada em feixes estreitos, entretanto, elas impõem o desafio do alinhamento angular rigoroso e contínuo com a espaçonave em movimento \cite{Ford2008}.
Isso ocorre porque o tempo de passagem de um satélite LEO pelo horizonte local é restrito a poucos minutos e qualquer desalinhamento no feixe principal ocasiona a perda de sincronismo e a degradação na qualidade do sinal recebido.

Para viabilizar o acompanhamento contínuo da trajetória orbital, empregam-se pedestais de rastreio automático.
Trata-se de mecanismos contendo de juntas motorizadas comumente dispostas em uma configuração bi-axial ortogonal de azimute e elevação.
Essa configuração torna os mecanismos capazes de orientar a antena ao longo de toda a abóbada celeste visível \cite{Bien2021}.
A precisão desse sistema depende da modelagem cinemática do mecanismo, do controle preciso de atuadores e da integração com algoritmos de estimação orbital em tempo real.

Em particular, a comunicação com satélites meteorológicos para recepção de dados no padrão HRPT (\textit{High Resolution Picture Transmission}) na faixa de $1,7\text{ GHz}$ demanda uma infraestrutura de recepção estável,
combinando antenas de polarização circular com sistemas de rastreio de baixo ruído e resposta rápida às mudanças de direção, necessários para garantir uma boa recepção do sinal.
Embora existam soluções comerciais de apontamento de antena, o custo elevado e a rigidez arquitetural limitam sua aplicação em estações acadêmicas e de radioamadorismo iniciante, justificando o desenvolvimento de protótipos embarcados customizados e de baixo custo.

Diante desse contexto, o objetivo deste trabalho é projetar, modelar e construir um protótipo funcional de um sistema de pedestal bi-axial com antena de geometria helicoidal para o rastreio automático de satélites LEO.
O projeto abrange a modelagem cinemática via parâmetros de Denavit-Hartenberg (D-H), a implementação embarcada de um filtro preditivo $\alpha-\beta-\gamma$ para suavização de trajetórias em microcontrolador dual-core (ESP32 via FreeRTOS), e a validação do sistema em um simulador para representar o comportamento real do mecanismo.

O presente relatório está estruturado da seguinte forma: o Capítulo~\ref{cap:obj} detalha os objetivos gerais e específicos, especificando cada requisito necessário para o sistema ser capaz de rastrear satélites; o Capítulo~\ref{cap:teo} fundamenta a teoria de rastreio, cinemática e filtragem preditiva nos quais o projeto foi fundamentado;
o Capítulo~\ref{cap:metod} descreve a metodologia de hardware, software e integração do sistema, passando pela escolha do software de rastreio, seleção dos componentes mecânicos e eletrônicos, implementação do filtro, arquitetura do software embarcado na placa, e o método de controle dos motores;
o Capítulo~\ref{cap:resultados} discute os resultados obtidos com o protótipo, explicitando as soluções obtidas para a montagem da antena, os métodos utilizados para a calibração dos parâmetros do filtro, as telas de interface do protótipo, e o funcionamento do modelo no simulador de física;
e, por fim, no Capítulo~\ref{cap:Conclusão} apresentam-se as conclusões e propostas de trabalhos futuros.
